\chapter{Disclosure risk as the outcome of a declared attack}\label{ch:scenario}

The minutes of the Bank's release committee contain one phrase more often than any other:
``every cell holds at least ten records, so no record carries a risk above one in ten.''
The chair wrote it, the committee approved it, and most readers of this book would approve it
too. It is the belief that a record's disclosure risk is the reciprocal of the number of
released records that share its key, so that a key class of ten means a one in ten chance.

The belief deserves its standing. The formula is old, it appears in every handbook, and it
answers a clear experiment: an adversary who knows a person's key, finds the released records
with that key, and picks one of them at random will be right one time in ten. The rule that
every class must hold at least $k$ records, known as $k$-anonymity, rests on it, and so do
most minimum cell counts in use at statistical institutions. If an adversary cannot tell the
records in a class apart, the formula is exact.

The fact that breaks it comes from a simulated register of 200,000 people with a 5 per cent
sample. A release that keeps at least ten records in every key class contains a record whose
own chance of being identified is 0.74, with a one-sided 95 per cent lower bound of 0.64.
% source: paper/generated/table_kanonymity.tex (k = 10 row: worst target 0.744, lower bound 0.638)
The adversary in that simulation knows one thing the key leaves out: a noisy reading of a
continuous variable that the release also publishes. That is enough to tell the ten records
apart.

This chapter asks one question: \emph{what does a disclosure risk figure measure, and under
what conditions does the reciprocal of the key frequency answer it?} The answer starts from
the objects every assessment involves, defines risk as the probability of one experiment, and
derives the familiar formula from that experiment under four assumptions. Stating the
assumptions shows where they fail, and two of them fail in the ordinary case.

\section{Population, sample, key and mechanism}

Let the finite population be $P = \{1,\dots,N\}$. The institution holds a
confidential sample, indexed by a random set $I \subseteq P$ of size $n$, and we
write $S_i = \ind{i \in I}$ for the inclusion indicator of unit $i$ and
$\pi_i = \Prob(S_i = 1)$ for its inclusion probability. A census is the case
$I = P$. Each unit has a confidential record $X_i$.

An adversary does not see all of $X_i$. It works through a key, a map $K_a$ from
records to a set of key values, where the subscript $a$ names the attack scenario.
The key collects the variables the adversary can observe about a person or firm
from outside: for a household survey perhaps age band, sex and district, for a
credit register perhaps sector, region and size class. For each key value $k$ the
population and sample frequencies are
\begin{equation}\label{eq:freq}
F_k = \sum_{i \in P} \ind{K_a(X_i) = k}, \qquad
f_k = \sum_{i \in P} S_i\, \ind{K_a(X_i) = k},
\end{equation}
so that $0 \le f_k \le F_k$. A record is sample unique when $f_k = 1$ and population
unique when $F_k = 1$.

The institution does not publish $X_I$ itself. It applies a protection mechanism,
which we write as a probability kernel $M_\theta$ indexed by a setting $\theta$, and
publishes a release $Y_\theta$ drawn from $M_\theta(\cdot \mid X_I)$. The notation
covers deterministic transformations, such as recoding or suppressing small cells,
as the case in which the kernel puts all its mass on one output. A minimum cell
count of five, a coarser age classification and a Laplace mechanism with
$\varepsilon = 1$ are all settings in this sense, and nothing so far requires a
setting to carry a formal guarantee.

Everything that follows depends on the randomness in the release, because a
disclosure risk is a probability over four sources at once: the sampling design,
the mechanism, the error in the adversary's own information, and any randomness in
the adversary's choice. An assessment that integrates over some of these and
conditions on others must say which, since the same file can give risk figures an
order of magnitude apart depending on that choice.

\section{One experiment}

Fix a target $t$ in the population. The adversary holds auxiliary information
$Z_t$ about the target, observes the release, and applies a linkage rule
$\delta_a$ that either names a released record as the target's or abstains.
Correct identification is the event that the target is in the file and the rule
names its record.

\begin{definition}[Scenario-conditional risk]\label{def:risk}
Given the background information $\mathcal{B}_a$ that scenario $a$ grants the
adversary, the unconditional and membership-conditional risks for target $t$ are
\begin{equation}\label{eq:tworisks}
R^{\mathrm{unc}}_a(t,\theta) = \Prob\{S_t = 1,\ \delta_a(Z_t,Y_\theta) = t \mid \mathcal{B}_a\},
\qquad
R^{\mathrm{in}}_a(t,\theta) = \Prob\{\delta_a(Z_t,Y_\theta) = t \mid S_t = 1,\ \mathcal{B}_a\}.
\end{equation}
\end{definition}

The two differ by the chance that the target is in the file at all. If a correct
link is impossible when the target is absent,
\begin{equation}\label{eq:factor}
R^{\mathrm{unc}}_a(t,\theta) = \pi_t\, R^{\mathrm{in}}_a(t,\theta).
\end{equation}
This is the multiplication rule of probability and involves no independence
assumption. It has practical weight because a sampling fraction is often applied
or omitted without comment, and with a 5 per cent sample that silent choice moves
the reported risk by a factor of twenty.

\section{Where the reciprocal of the key frequency comes from}

The most familiar disclosure measure assigns a record with key $k$ the risk
$1/f_k$: if three released records share the target's key, the adversary has a one
in three chance. The formula is usually presented as self-evident. It is in fact
the answer to the experiment above under four assumptions, which are worth stating
one by one.

\begin{assumption}\label{as:key}
The released key is unperturbed: applied to the release, the key map gives the
same values as applied to the confidential sample.
\end{assumption}
\begin{assumption}\label{as:obs}
The adversary observes the target's key without error.
\end{assumption}
\begin{assumption}\label{as:exch}
All records that share the target's key are exchangeable in the adversary's
information: nothing the adversary knows favours one of them over another.
\end{assumption}
\begin{assumption}\label{as:unif}
The adversary chooses uniformly among the records compatible with what it knows.
\end{assumption}

\begin{proposition}[Reciprocal frequency]\label{prop:recip}
Under Assumptions~\ref{as:key} to~\ref{as:unif}, for a target in the file whose key
is $k$,
\begin{equation}\label{eq:recip}
\Prob\{\delta_a = t \mid S_t = 1,\ K_a(X_t) = k,\ f_k\} = \frac{1}{f_k}.
\end{equation}
\end{proposition}

The result is old \citep{duncan1989,skinner2002}, and the value of restating it
lies in the list of conditions. Assumption~\ref{as:exch} fails whenever the
adversary holds a variable, correlated with something the release publishes, that
distinguishes candidates within a key class. In financial and health data that is
the normal case: a lender knows roughly how large its own loan to a firm is, and a
neighbour knows roughly how old someone is. Assumption~\ref{as:obs} fails whenever
the adversary's own source contains reporting error, which is also the normal case.
When either fails, the right object is no longer a frequency but a posterior
probability.

At the Bank this is where the chair's phrase runs into trouble. The credit register publishes
each loan's exposure beside its sector, region and size class. Pieter, a credit analyst at
one of the reporting lenders, knows the size of every loan his bank has made, give or take a
rounding. When he looks for one of his own borrowers in a class of ten, he does not choose at
random. He chooses the record whose exposure matches his books, and Assumption~\ref{as:exch}
no longer holds.

\section{When the assumptions fail}

Suppose the adversary's information and the released records are compared through
a likelihood, in the sense of record linkage under error \citep{fellegi1969}. Let
$L_{tj}$ be the likelihood of what the adversary observes under the hypothesis that
released record $j$ belongs to target $t$, let $q_{tj}$ be the prior weight on that
hypothesis, and let $j = 0$ stand for the hypothesis that the target is not in the
file. Bayes' rule gives the posterior
\begin{equation}\label{eq:post}
p_{tj} = \frac{q_{tj}L_{tj}}{q_{t0}L_{t0} + \sum_{\ell \in I} q_{t\ell}L_{t\ell}},
\end{equation}
and an adversary that names the most probable candidate succeeds with probability
\begin{equation}\label{eq:maxpost}
R^{\mathrm{in}}_a(t,\theta) = \E\Big[\max_{j \in I} p_{tj} \ \Big|\ S_t = 1,\ \mathcal{B}_a\Big].
\end{equation}
When the compatible records have equal likelihood and the no-match prior is zero,
\eqref{eq:post} and \eqref{eq:maxpost} reduce to the reciprocal frequency. In every
other case they can lie above or below it.

Two consequences follow, and both cut against common practice. First, measurement
error cannot be assumed to reduce risk. It changes the likelihood ratios, and it
may spread the posterior across candidates or concentrate it on one. In the
simulation study of Part~V, adding reporting error to the key lowers the
adversary's success from 0.466 to 0.381, while giving the adversary one noisy
reading of a continuous variable brings it back to 0.464. The reciprocal-frequency
figure was close to the truth in that case by cancellation, not by argument.

Second, a release can satisfy a minimum cell count on its declared key and still
expose individual records. If the release publishes a continuous variable, such as
an exposure or an income, and the adversary holds a noisy reading of it, the
posterior concentrates on the record whose value is nearest the reading, whatever
the size of the key class.

\section{$k$-anonymity is a threshold statement}

The rule that every released key class should contain at least $k$ records,
$k$-anonymity in the sense of \citet{sweeney2002}, is the most widely used
protection criterion in this field. The experiment above shows exactly what it
delivers.

\begin{proposition}[$k$-anonymity as a decision rule]\label{prop:kanon}
Under Assumptions~\ref{as:key} to~\ref{as:unif}, and for a tolerance
$\tau \in (0,1]$, the largest prosecutor risk in the release is at most $\tau$ if
and only if every occupied key class has at least $\lceil 1/\tau \rceil$ members:
\begin{equation}\label{eq:kanon}
\sup_{t \in I} \frac{1}{f_{K_a(X_t)}} \le \tau
\iff \min_{k:\,f_k>0} f_k \ge \lceil 1/\tau \rceil .
\end{equation}
If Assumption~\ref{as:exch} fails, for instance because the release publishes a
continuous variable of which the adversary holds a noisy reading, the implication
from right to left fails: there are releases that are $k$-anonymous on the declared
key in which some target's risk exceeds $1/k$.
\end{proposition}

Read this way, $k$-anonymity is a special case of the decision rule developed in
Part~III. It is the case in which the risk model is exact exchangeable matching,
the estimate carries no uncertainty, and the tolerance is the reciprocal of an
integer. Each of these simplifications is a place where an institution can be
wrong, and only the third is a policy choice. The failure in the second half of
the proposition is not hypothetical. On a simulated register of 200,000 units with
a 5 per cent sample, a release that is 10-anonymous on its declared key contains a
record whose own re-identification probability is 0.74, with a one-sided 95 per
cent lower bound of 0.64. Raising the minimum class size on that key does not
establish a bound for the enriched attack, because the variable that separates the
records is not part of the key.

The refinements of $k$-anonymity follow the same pattern. The $\ell$-diversity
criterion \citep{machanavajjhala2007} addresses the case in which all records in a
key class share a sensitive value, and $t$-closeness \citep{li2007} limits how far
the distribution of a sensitive value within a class may differ from its overall
distribution. Both protect against attribute disclosure, which can occur without
any record being identified. They change the risk model rather than the decision
rule, and they cannot be read as interchangeable tolerances for correct linkage.

Naledi ran Pieter's attack on the register before the next meeting. Against an adversary who
reads exposure with small error, the minimum cell count that the committee had treated as a
one in ten guarantee bounded nothing at all: Chapter~\ref{ch:credit} reports that no floor up
to 300 brings his success below the committee's tolerance. The chair's phrase was true of a
different adversary, one who knows only the key.

\section{What a risk figure must carry}

The argument of this chapter reduces to a requirement. A reported risk figure
should specify its scenario: who the adversary is, what it knows, which targets it
attacks, and which conditioning event the figure assumes, in particular whether the
target is known to be in the file. Chapter~3 shows that three measures commonly
reported under one label, all computed on the same simulated release, give 0.466,
0.064 and 0.023. Each is correct for its own experiment, and only a declared
scenario says which experiment the institution meant.

\section{The answer}

The question was what a disclosure risk figure measures and when the reciprocal of the key
frequency answers it. A risk figure measures the probability that a declared adversary, with
declared knowledge, names the right record in one defined experiment, and it means nothing
until that experiment is stated. The reciprocal of the key frequency is the exact answer when
the released key is unperturbed, the adversary observes it without error, nothing it knows
distinguishes the records in a class, and it chooses among them at random. When the adversary
holds a variable the release also publishes, or observes the key with error, the answer is a
posterior probability that can lie well above or below $1/f_k$, and a minimum cell count on
the key alone no longer bounds it. The chair's minutes should have read: ``no record carries a
risk above one in ten against an adversary who knows only the key.''

\begin{inpractice}
For every risk figure in a release assessment, record the adversary's auxiliary
variables and how precisely it observes them, whether the target is assumed to be
in the file, and whether the release publishes any variable outside the key that the
adversary might also know. If it does, a minimum cell count on the key alone does
not bound the risk, and the assessment needs a simulated attack that uses that
variable.

\smallskip\noindent\textit{In the software.} The measures of this chapter are in
\texttt{privacy\_risk.risk} (\texttt{prosecutor\_risk},
\texttt{marketer\_sample\_risk}, \texttt{marketer\_unconditional\_risk}), and the
posterior linkage of \eqref{eq:post} and \eqref{eq:maxpost} is in
\texttt{privacy\_risk.linkage} (\texttt{posterior\_linkage},
\texttt{maximum\_posterior\_success}, and \texttt{attack} for the simulated
version).
\end{inpractice}

\section{Proofs}

\begin{proof}[Proof of Proposition~\ref{prop:recip}]
Condition on $S_t = 1$ and $K_a(X_t) = k$. By Assumption~\ref{as:key} every
released record carries its confidential key, and by Assumption~\ref{as:obs} the
adversary observes $k$. The compatible set is therefore
$\{j \in I : K_a(X_j) = k\}$, which has $f_k$ members by \eqref{eq:freq} and
contains $t$. By Assumption~\ref{as:exch} the adversary's information is invariant
to permutations of that set, and by Assumption~\ref{as:unif} it selects each member
with probability $1/f_k$. Exactly one member is the target, so the probability of
a correct choice is $\sum_{j} f_k^{-1}\ind{j = t} = 1/f_k$.
\end{proof}

\begin{proof}[Proof of Proposition~\ref{prop:kanon}]
By Proposition~\ref{prop:recip}, the largest prosecutor risk in the release is
$1/\min_{k:\,f_k>0} f_k$. It is at most $\tau$ if and only if
$\min_k f_k \ge 1/\tau$, and since frequencies are integers this holds if and only
if $\min_k f_k \ge \lceil 1/\tau \rceil$, which is $k$-anonymity with
$k = \lceil 1/\tau \rceil$.

For the failure claim, drop Assumption~\ref{as:exch}. Let the release publish a
confidential variable $y$ beside the key, let a key class have exactly $k$ members
with distinct values of $y$, and let the adversary observe
$z_t = y_t + \sigma\xi$ with $\xi$ standard normal. Under the Gaussian likelihood
the posterior \eqref{eq:post} is largest for the member whose $y$ is nearest $z_t$.
Let $d > 0$ be the smallest gap between two values of $y$ in the class. The
maximum-posterior rule is correct whenever $|\sigma\xi| < d/2$, an event whose
probability tends to one as $\sigma \to 0$. Hence for all sufficiently small
$\sigma$ the target's risk exceeds $1/k$, while the release is $k$-anonymous on the
key for every $\sigma$.
\end{proof}
